\chapter{Hardware}

\section{Circuit of the TransistorTester}
\label{sec:hardware}
The circuit of the TransistorTester in figure \ref{fig:ttester} is based on the circuit of
 Markus F. released in Abb.~1 of AVR-Transistortester report \cite{Frejek}.
Changed or moved parts are marked with green color, optional parts are marked with red color.

Some changes are made because the electronical power switch make problems in some implementations.
Therefore the resistor R7 is reduced to \(3.3k\Omega\). The capacitor C2 is reduced to
10nF and R8 is moved so that the PD6 output does not try to switch the C2 capacitor directly.
 Additional blocking capacitors are added and should be placed
near the power connection of the Atmega and near the Voltage regulator.

Because the PD7 input and PC6 (RESET) are the only pins, where pull up resistors where
needed, one extra \(27k\Omega\) resistor is added to the PD7 (pin~13) input. With this modification
the software can disable all internal pull up resistors of the ATmega.

The additional crystal with its 22pF capacitors are optional added. 
The accuracy of a crystal has the benefit of more stable time measurement for getting the
capacitor values.

New software version can use a voltage scale switch of the ADC. The speed of switching is reduced
by the external capacitor C1 at the AREF (21) pin of the ATmega. To avoid slowing down the
measurement speed more than necessary, the value of this capacitor should be reduced to 1nF.
Removing of the capacitor C1 is also possible.
For adapting the software to the actual circuit take a look to the Makefile options in the
configuring chapter \ref{sec:config}.

Some different versions of R11 / R12 resistor combinations circulates
in the internet. I~have adapted my software to the original of Markus F. \cite{Frejek} with \(10k\Omega\) and \(3.3k\Omega\).

The additional 2.5V precision voltage reference connected at pin PC4 (ADC4) can be used to
check and calibrate the VCC voltage, but is not required.
A optional ISP connector has been added to easier load new software versions to the tester.


\begin{figure}[H]
\centering
\includegraphics[width=18cm]{../FIG/ttester.eps}
\caption{New circuit of TransistorTester}
\label{fig:ttester}
\end{figure}
The software can follow to another pin assignment of port D for a simpler connection of the
LCD display. 
The following table \ref{tab:grid-change} shows the modified assignments.

\begin{table}[H]
  \begin{center}
    \begin{tabular}{| c | c | c |}
    \hline
       Signal & circuit diagram & strip grid board version\\
    \hline
    pushbutton input  &  PD7   &  PD0 \\
    LCD-RS    &  PD4      & PD7 \\
    LCD-E     &  PD5   & PD5 \\
    LCD-D4    &  PD0   & PD4 \\
    LCD-D5    &  PD1   & PD3 \\
    LCD-D6    &  PD2   & PD2 \\
    LCD-D7    &  PD3   & PD1 \\
    \hline
    \end{tabular}
  \end{center}
  \caption{Changes for strip grid board}
  \label{tab:grid-change}
\end{table}

For better protection of the ATmega inputs the additional circuit~ \ref{fig:relay_addon} 
can be integrated. The de-energized contacts of the relay protect the ATmega without power.
The contacts will be opened by software only for measurement.

\begin{figure}[H]
\centering
\includegraphics[width=7cm]{../FIG/relay_addon.eps}
\caption{Additional protection of the ATmega inputs}
\label{fig:relay_addon}
\end{figure}



\section{Hints for building the TransistorTester}
Every LCD-display with at least 2x16 character and a HD44780 compatible controller
can be used with the TransistorTester. You should respect the current needed for
illumination, some LCD need lower current than others.
I had tried OLED type displays, but this type cause interference with measurements
of the ATmega and are {\bf not} recommended. Also loading of special characters
for displaying the resistor symbol has caused problems with the OLED.

The resistors R1 to R6 are critical for measurements and this \(680\Omega\) and
\(470k\Omega\) resistors should be measurement type resistors (tolerance of 0.1\%) to
get the full accuracy.
You should use a precision socket for the ATmega microcontroller to enable
the replacement of the microcontroller.
The microcontroller ATmega8, ATmega168 and ATmega328 can be used.
Recommended is a ATmega168 or ATmega328, if you wish to use all features.

Anyway you should assemble all parts to printed board without the microcontroller.
A up-to-date low voltage drop regulator like MCP1702-5002 is recommended as IC2, because it
need only \(2\mu A\) of standby current and can still deliver 5V, if your input
voltage is only 5.4V. But this part is not pin compatible to well known 78L05 with TO92 body!

After checking, that all needed parts are at the correct place, you should
first connect the battery or power supply to the printed board without LCD-display
and microcontroller. You should check the voltage at the power pins of the
microcontroller and LCD-display terminal during the Test key is pressed.
The voltage should disappear, if you release the Test key.
If the voltage had correct polarity and value,
you should disconnect the power and assemble the microcontroller with correct
alignment. Be careful and make shure, that all pins of the microcontroller
are in the socket holes.
Now you can also connect the LCD. Check if power pins of the LCD has the right connection to
GND and VCC of your board.

If you are shure that everything is all right, reconnect the power. 
If you have already programmed the ATmega, you can press the Test button.
By pressing the Test key, the background light of the LCD should switch on.
If you release the Test button, the LED should illuminate weak.
Notice, that the software for the microcontroller must be compiled for the
correct processor type. A program for the ATmega8 does not run on a ATmega168!

\section{Changeover for tester versions designed by Markus F.}
\label{sec:change_markus}
\begin{description}

\item[Voltage control]
If the problem exist, the tester will shut down immediately with every switch on.
With imy suggested setting of the fuses (Makefile) the voltage control of the different
ATmega versions is switched to 4V (brown out level).
This may be the reason why the tester makes trouble with the power on sequence.
The Pin PD6 tries to switch the 100nF capacitor C2 to VCC level causing a voltage
breakdown of the VCC voltage (5V).
The capacitor C2 can be reduced to \textless 10nF without problems.
If possible, the direct connection of PD6 should be replaced by a resistor \textgreater \(220 \Omega\).
\item[Improvement of power on circuit]
Often this problem is the reason, if the tester starts with the button hold pressed, but switch off
directly by releasing. The problem is enforced by a high current background light for the LCD.
The resistor R7 to the base of the PNP transistor T3 was optimized with the value \(27k \Omega\) 
too much to save power consuming.
To improve the switching with lower battery voltage or lower current amplification factor of
the PNP transistor T3, you should reduce the resistance to \(3.3k \Omega\).
\item[Additional pull-up resistor at PD7]
The missing pull-up resistor results to a switch off of the tester with the message ''Timeout''
after a short display time.
The software is configured with the option PULLUP\_DISABLE, that all internal pull-up
resistors are switched off. For that reason the voltage of pin PD7 is not definded,
if the level is not switched by the push button or transistor T2 to GND.
One external pull-up resistor of \(27k \Omega\) to VCC avoid this error.
\item[Capacitor C1 at the AREF pin]
Many designs use a 100nF capacitor at the AREF pin, like the design of Markus F. too.
As long as the reference voltage of the ADC is never changed, this is a good solution.
The software of the TransistorTester for the ATmega168/328 uses a automatic selection
of the internal 1.1 V reference voltage of the ADC, if the input voltage is below 1V.
With this solution a better resolution of the ADC can be reached for little input voltages.
Unfortunately the switching from 5V to 1.1V reference is very slowly. A additional
wait time of 10ms must be respected for this reason.
With changing the capacity value to 1nF this wait time can be reduced significant.
I have not noticed any degration of measurement quality with this change.
Even a removing of the capacitor has no significant change of measurement results.
If you prefer to leve the capacitor unchanged, you can remove the option NO\_AREF\_CAP
in the Makefile to activate longer wait times in the program.
\item[Expanding of a 8MHz crystal]
With some skill you can expand a 8MHz crystal to the backside of the printed board
directly to the pins PB6 and PB7 (pin 9 and pin 10).
My own expansion was done without the both 22pF capacitors.
This solution has operated well with all tested ATmega.
But it is not required to use a crystal. You can still use the 8MHz RC oszillator
by setting the fuses to get the better resolution of time constant measuring (capacity value).
\item[Smoothing of the operating voltage]
The original circuit of Markus F. shows only one 100nF capacitor to block the VCC voltage.
This is clearly too little smoothing. You should at least use one 100nF near the ATmega power pins
and one near the voltage regulator. The input of the voltage regulator should be
blocked with a 100nF too.
Additional \(10\mu F\) capacitors (electrolytic or ceramic) at the input and
output of the voltage regulator can stable the voltage level.
Ceramic \(10\mu F\) capacitos with SMD mounting form are easier to use for backfitting
and have usually a lower ESR value. 
\item[Selection of the ATmega processor]
The using of the base function of the tester is still possible with a ATmega8.
The flash memory of that device is used near 100\% .
Because the ATmega168 or ATmeg328 processors are pin-compatible to the ATmega8,
I can recommend the replacement.
Actually the price for ATmega328 is so cheap, that there is no reason to take
a ATmega168 type.
With a ATmega168/328 you get the following advantages:
Self test with automatic calibration.\\
Improvement of measurement quality by automatic switching of ADC scale.\\
Measurement of inductors with resistance  below \(2100 \Omega\).\\
Measurement of ESR value of capacitors with value of above  \(2 \mu F\).\\
Using of pin PC4 as serial output.\\
\item[Missing precision voltage reference]
Usually the software should detect the missing voltage reference with the unconnected pin PC4.
In this case no VCC=x.xV message should appear in row 2 of the LCD on power on.
If this message appear without the reference, you should connect a \(2.2k \Omega\) resistor
to the PC4 input and VCC.


\end{description}

\section{Chinese clones}
As I know, the tester is rebuild in China in two versions.
The first model is rebuild from the first design of Markus F. without the ISP port.
The assembled ATmega8 is placed in a socket, so you can replace it with a ATmega168 or ATmega328.
For this version you should consider all the hints of section \ref{sec:change_markus}.
Additional \(100nF\) ceramic cpacitors should be connected near by the VCC-GND and AVCC-GND pins of
the ATmega for better stabilization of the power voltage.
In addition you should notice, that if you expand the board with the additional 8 MHz crystal,
your external ISP programmer must have a external clock for programming.\\

The second version of rebuilded tester is build with SMD components. Also the fix installed ATmega168
is a SMD type with 32TQFP body.
Fortunately on the board is a 10-pole ISP connector provided for the programming.
I have analysed the board version ''2.1 2012/11/06''. One error is the assembly of the part ''D1'',
which should be a precision 2.5V voltage reference. Assembled is only a zener diode.
This part should be removed. You can mount a LM4040AIZ2.5 or LT1004CZ-2.5 precision voltage reference
at this place. A missing voltage reference is noticed by the software, so that you must not install
the voltage reference.
My exemplar was delivered with software version 1.02k. The 10-pole ISP plug was not assembled and I must
install a jumper from ISP pin 6 to ISP pin 10. My programmer expect a GND connection at pin 10, but the
board has GND level only on pin 4 and pin 6 of the ISP.
The label of the ATmega168 was rub away and there was no documentation delivered with the part.
The lock fuses of the ATmega were set, so no readout was possible.
But I could install the software version 1.05k without any problems.
Another user has problems with the same software version 1.05k. This user has the chinese board ''2.2 2012/11/26''.
The software runs only without problems, if a additional \(100nF\) keramic capacitor was placed between
the pin 18-AVCC and 21-GND near by the ATmega.
The software 1.05k uses the sleep state of the ATmega for waiting time. For this reason the current alternates
often and the voltage regulator is stressed more.
Further I have noticed, that the VCC voltage is blocked with a \(100nF\) ceramic capacitor and with a
\(220\mu F\) electrolytic capacitor nearby the 78L05 voltage regulator.
The 9V supply voltage is blocked with the same capacitors, but not at the input of the regulator but
at the emitter of the PNP transistor (parallel with the battery). 
The printed circuit board track from the ATmega168 to the test port is very thin, so that a resistance
of \(100m \Omega\) could be measured for one path. This will be the reason for measuring a resistance
of \(0.3 \Omega\) for two direct connected pins.
The ESR measuring can usually consider this by zero compensation.
The current version of software does not respect this offset for measuring of resistors with low resistance.


\section{Programming of the microcontroller}
I release the software for the microcontroller with source code.
The developement is done with Linux operationg system (Ubuntu) and
is controlled with a Makefile. The Makefile makes shure, that your
software will be compiled with the prior selected Makefile options. Some constellations
are precompiled with the source. Please take a look to the ReadMe.txt file
in the directory Sourcecode/default and to the chapter~\ref{sec:config}.
The result of compilation have the extensions .hex and .eep .
Usually the names will be TransistorTester.hex and TransistorTester.eep .
The .hex file contains the data for the program memory (flash) of the ATmega processor.
The .eep file contains the data for the EEprom memory of the ATmega. Both data files
must be loaded to the correct memory.

Additionally the operating state of the
ATmega processor must be programmed with the ``fuses''.
If you can use my Makefile and additionally the program avrdude \cite{avrdude}, you need no exact
knowledge of the details about the fuses. You have only to type ``make fuses'' if you
have no crystal or ``make fuses-crystal'' if you have installed the 8MHz crystal to your printed board.
With the ATmega168 series of the microcontroller you can also use ``make fuses-crystal-lp'' to use
a crytal with the low power mode.
Never choose the crystal mode of clock generation, if you don't have installed
the 8MHz crystal. If you are not shure with the fuses, leave them as default
set by manufactor and first bring the the tester to operation in this mode.
Maybe your program runs too slow, if you use program data compiled for
8MHz operation, but you can correct this later! But a wrong set of fuses may inhibit
later ISP-programming.
If you use the Windows operating system, the easiest way to get a correct programmed
ATmega is to use the WinAVR package \cite{winavr1},\cite{winavr2}.
With my patch \cite{winavr3} you can also set the fuses by using the Makefile.
Of course the avrdude program must support your programmer and the configuration
in the Makefile must match to your environment.


\section{Troubleshooting}
In most cases of problems you will miss the text output to the LCD-display.
At first you should check, if the LED was illuminated weak, if you release
the Test button. 
\begin{description}

\item[Power does not switch on.]
If the LED is without light and the VCC power has correct
5V voltage during holding the Test button, the microcontroller does not switch the power
correctly. The microcontroller should hold the power by switching the
PD6 output to 5V, which is usually done as one of the first actions.
If you hold the Test key pressed, the power is switched on anyway.
So you can check the value of VCC power and additionally the voltage value
of the PD6 output, if you hold the key pressed.
If VCC voltage has correct value (5V), but PD6 voltage is
below 4V, your microcontroller does not start the program. In this case
you should check if the microcontroller flash has been loaded with proper data for your
installed type and if ATmega is correctly configured with the fuses.
If your ATmega put the PD6 output to 5V and the power does not stay if you
release the Test key, it is more difficult to find the reason.
First you can shorten the LED and try again. If your Tester now starts,
your LED may be faulty or mounted with wrong polarity. If this is not
the reason, the current amplification factor of your T3 transistor (BC557C)
is insufficient. The current to the base of T3 is lower in the microcontroller
state as in the ``key pressed'' state.

\item[Nothing is readable on the LCD display]
Check the voltage at the contrast pin at the LCD display (pin 3). Adjust to
correct value specified in the data sheet of your display and optimize by viewing.
If you have a high temperature display type, you must provide a negative contrast voltage
for operation. In this case you can use the ICL~7660 device for generating
a negative voltage from positive 5V.

If there is no output readable on the LCD and the background light is on,
you should disconnect the power and check all four data plus the two control signal connections.
If all connection are well, the only reason I see is a uncorrect timing of
control signals. This can be caused by a slower LCD controller than expected by
the software or the ATmega software runs at wrong clock speed. Please check for which
clock speed your programming data was compiled  and if the fuses of the
ATmega are correct set to that speed. You find the clock parameter in the corresponding
Makefile.
If the tester is build without the switch off electronic, you can test with
a LED connected to the test pins, if the program operates normally.
If the LED flickers, the program operates well. The missing text on the
LCD must be caused by wrong connection or timing.


\item[Something but not all is readable on the LCD display]
Check if the .eep data are loaded to the EEprom memory of ATmega.
If all data are loaded correctly, you should check the clock speed of your
programming data (Makefile) and ATmega processor settings (fuses).

\item[Measurement is slow and Capacitors are measured about 8 times too small]
You run software compiled for 8MHz clock at real clock speed of 1MHz.
Please set the fuses of the ATmega correctly.

\item[Measurement has strangely values]
Check if your programmer is still connected to the ISP-plug.
The ISP interface should be disconnected for measuring.
Very often the reason of wrong measurements is the use of software compiled with
the AUTOSCALE\_ADC option and with the option NO\_REF\_CAP, but the capacitor
at the AREF pin has still a value of 100nF.
Wrong assembly of components or remaining soft solder flux can disturb the 
measurements too. Please check with the selftest function of your TransistorTester software
if possible. For the details see Chapter \ref{sec:selftest}.

Otherwise inspect your board visually and check the resistor values
with a ohmmeter. You can use the pins of the ATmega for this check, for example
to check the R1 you can measure between pin 23 and pin 14. Take a look at the
circuit diagram \ref{fig:ttester} for details. There is no need to
remove the microcontroller, only battery or power supply should be removed before.


\end{description}
